\documentclass{article}
\usepackage{graphicx} % Required for inserting images
\usepackage{amsmath}
\usepackage{amssymb}
\usepackage{algorithm}
\usepackage{algpseudocode}

\title{Sutton and Barton Exercises Chapter 12}
\author{danieltocco0 }
\date{April 2026}

\begin{document}

\maketitle

\section{Chapter 12}


\textbf{Exercise 12.1} \quad Just as the return can be written recursively in terms of the first reward and 
itself one-step later (3.9), so can the $\lambda$-return. Derive the analogous recursive relationship 
from (12.2) and (12.1).
\\

\begin{align*}
G_t^\lambda 
&= (1-\lambda)\sum_{n=1}^{\infty} \lambda^{n-1} G_t^{(n)} \\
&= (1-\lambda)\sum_{n=1}^{\infty} \lambda^{n-1} \left[ R_{t+1} + \gamma G_{t+1}^{(n-1)} \right] \\
&= (1-\lambda)\sum_{n=1}^{\infty} \lambda^{n-1} R_{t+1}
+ (1-\lambda)\sum_{n=1}^{\infty} \lambda^{n-1} \gamma G_{t+1}^{(n-1)} \\
&= R_{t+1} + \gamma (1-\lambda)\sum_{n=1}^{\infty} \lambda^{n-1} G_{t+1}^{(n-1)} \\
&= R_{t+1} + \gamma (1-\lambda)\sum_{k=0}^{\infty} \lambda^k G_{t+1}^{(k)} \\
&= R_{t+1} + \gamma (1-\lambda)\left[ G_{t+1}^{(0)} + \sum_{k=1}^{\infty} \lambda^k G_{t+1}^{(k)} \right] \\
&= R_{t+1} + \gamma (1-\lambda) V(S_{t+1})
+ \gamma (1-\lambda)\sum_{k=1}^{\infty} \lambda^k G_{t+1}^{(k)} \\
&= R_{t+1} + \gamma (1-\lambda) V(S_{t+1})
+ \gamma \lambda \left[(1-\lambda)\sum_{k=1}^{\infty} \lambda^{k-1} G_{t+1}^{(k)} \right] \\
&= R_{t+1} + \gamma (1-\lambda) V(S_{t+1})
+ \gamma \lambda G_{t+1}^\lambda \\
&= R_{t+1} + \gamma \left[(1-\lambda) V(S_{t+1}) + \lambda G_{t+1}^\lambda \right].
\end{align*}
\\\\
\textbf{Exercise 12.2} \quad The parameter $\lambda$ characterizes how fast the exponential weighting in 
Figure 12.2 falls off, and thus how far into the future the $\lambda$-return algorithm looks in 
determining its update. But a rate factor such as $\lambda$ is sometimes an awkward way of 
characterizing the speed of the decay. For some purposes it is better to specify a time 
constant, or half-life. What is the equation relating $\lambda$ and the half-life $\tau_\lambda$, the time by 
which the weighting sequence will have fallen to half of its initial value?
\\
\\

\textbf{Solution:}

\text{The weight at step } $n$ \text{ is } $(1-\lambda)\lambda^{n-1}$\text{. The initial value at } $n=1$ \text{ is } $(1-\lambda)$.
\text{ We want } $\tau_\lambda$ \text{ such that the weight has fallen to half:}

\[
(1-\lambda)\lambda^{\tau_\lambda - 1} = \frac{1}{2}(1-\lambda)
\]

\text{Cancel } $(1-\lambda)$\text{:}
\[
\lambda^{\tau_\lambda - 1} = \frac{1}{2}
\]

\text{Take the natural log of both sides:}
\[
(\tau_\lambda - 1)\ln\lambda = -\ln 2
\]

\text{Solve for } $\tau_\lambda$\text{:}
\[
\boxed{\tau_\lambda = 1 - \frac{\ln 2}{\ln \lambda}}
\]
\\\\
\textbf{Exercise 12.3} \quad Some insight into how $\text{TD}(\lambda)$ can closely approximate the 
off-line $\lambda$-return algorithm can be gained by seeing that the latter's error term 
(in brackets in (12.4)) can be written as the sum of TD errors (12.6) for a single fixed 
$\mathbf{w}$. Show this, following the pattern of (6.6), and using the recursive relationship 
for the $\lambda$-return you obtained in Exercise 12.1.
\\
\\

\textbf{Solution:}

\text{Recall from (6.6) that for fixed } $\mathbf{w}$\text{:}
\[
G_t - \hat{v}(S_t, \mathbf{w}) = \sum_{k=t}^{T-1}\gamma^{k-t}\delta_k
\]

\text{We want the analogous result for the } $\lambda$\text{-return error. From Exercise 12.1:}
\[
G_t^\lambda = R_{t+1} + \gamma\left[(1-\lambda)\hat{v}(S_{t+1},\mathbf{w}) + \lambda G_{t+1}^\lambda\right]
\]

\text{Subtract } $\hat{v}(S_t, \mathbf{w})$ \text{ from both sides:}
\[
G_t^\lambda - \hat{v}(S_t,\mathbf{w}) = R_{t+1} + \gamma(1-\lambda)\hat{v}(S_{t+1},\mathbf{w}) 
+ \gamma\lambda G_{t+1}^\lambda - \hat{v}(S_t,\mathbf{w})
\]

\text{Add and subtract } $\gamma\hat{v}(S_{t+1},\mathbf{w})$\text{:}
\[
= \underbrace{R_{t+1} + \gamma\hat{v}(S_{t+1},\mathbf{w}) - \hat{v}(S_t,\mathbf{w})}_{\delta_t} 
+ \gamma\lambda\bigl(G_{t+1}^\lambda - \hat{v}(S_{t+1},\mathbf{w})\bigr)
\]

\text{Expanding recursively:}
\[
= \delta_t + \gamma\lambda\delta_{t+1} + (\gamma\lambda)^2\delta_{t+2} + \cdots
\]

\[
\boxed{G_t^\lambda - \hat{v}(S_t,\mathbf{w}) = \sum_{k=t}^{T-1}(\gamma\lambda)^{k-t}\delta_k}
\]
\\\\
\textbf{Exercise 12.4} \quad Use your result from the preceding exercise to show that, if the 
weight updates over an episode were computed on each step but not actually used to change 
the weights ($\mathbf{w}$ remained fixed), then the sum of $\text{TD}(\lambda)$'s weight updates would be 
the same as the sum of the off-line $\lambda$-return algorithm's updates.

\text{As } $\lambda \to 0$\text{, } $\tau_\lambda \to 1$ \text{ (weights die after one step).}
\text{ As } $\lambda \to 1$\text{, } $\tau_\lambda \to \infty$ \text{ (weights decay very slowly, approaching Monte Carlo).}
\\
\\

\textbf{Solution:}

\text{The } $\text{TD}(\lambda)$ \text{ update at each step with eligibility trace } 
$\mathbf{z}_t = \sum_{k=0}^{t}(\gamma\lambda)^{t-k}\nabla\hat{v}(S_k,\mathbf{w})$ \text{ is:}
\[
\Delta\mathbf{w}_t = \alpha\delta_t\mathbf{z}_t
\]

\text{The sum of all updates over the episode is:}
\[
\sum_{t=0}^{T-1}\alpha\delta_t\mathbf{z}_t 
= \alpha\sum_{t=0}^{T-1}\delta_t\sum_{k=0}^{t}(\gamma\lambda)^{t-k}\nabla\hat{v}(S_k,\mathbf{w})
\]

\text{Swap the order of summation:}
\[
= \alpha\sum_{k=0}^{T-1}\nabla\hat{v}(S_k,\mathbf{w})\sum_{t=k}^{T-1}(\gamma\lambda)^{t-k}\delta_t
\]

\text{Apply the result from Exercise 12.3:}
\[
\sum_{t=k}^{T-1}(\gamma\lambda)^{t-k}\delta_t = G_k^\lambda - \hat{v}(S_k,\mathbf{w})
\]

\text{Substituting:}
\[
\boxed{\sum_{t=0}^{T-1}\alpha\delta_t\mathbf{z}_t 
= \alpha\sum_{k=0}^{T-1}\nabla\hat{v}(S_k,\mathbf{w})\bigl(G_k^\lambda - \hat{v}(S_k,\mathbf{w})\bigr)}
\]

\text{Which is exactly the sum of the offline } $\lambda$\text{-return algorithm's updates — one}
\text{gradient step per timestep toward the } $\lambda$\text{-return target.}
\\\\
\textbf{Exercise 12.5} \quad Several times in this book we have established that returns 
can be written as sums of TD errors if the value function is held constant. Why is (12.10) 
another instance of this? Prove (12.10).
\\
\\

\textbf{Explanation:}

\text{Equation (12.10) is another instance of the sum-of-TD-errors pattern because each } 
$\delta_i'$ \text{ treats the weights as locally fixed for that one step — } $\mathbf{w}_{i-1}$ 
\text{ for the current state and } $\mathbf{w}_i$ \text{ for the next state. The telescoping }
\text{cancellation still holds step-by-step, just with a slightly different } $\delta'$ 
\text{ that accounts for the weight update between steps.}

\textbf{Proof of (12.10):}

\text{We want to show:}
\[
G_{t:t+k}^\lambda = \hat{v}(S_t, \mathbf{w}_{t-1}) + \sum_{i=t}^{t+k-1}(\gamma\lambda)^{i-t}\delta_i'
\]

\text{where } $\delta_i' = R_{i+1} + \gamma\hat{v}(S_{i+1}, \mathbf{w}_i) - \hat{v}(S_i, \mathbf{w}_{i-1})$.

\text{Expand the sum by writing out each term:}
\[
\sum_{i=t}^{t+k-1}(\gamma\lambda)^{i-t}\delta_i' 
= \delta_t' + \gamma\lambda\delta_{t+1}' + (\gamma\lambda)^2\delta_{t+2}' + \cdots + (\gamma\lambda)^{k-1}\delta_{t+k-1}'
\]

\text{Substituting } $\delta_i' = R_{i+1} + \gamma\hat{v}(S_{i+1},\mathbf{w}_i) - \hat{v}(S_i,\mathbf{w}_{i-1})$\text{:}

\[
= \bigl[R_{t+1} + \gamma\hat{v}(S_{t+1},\mathbf{w}_t) - \hat{v}(S_t,\mathbf{w}_{t-1})\bigr]
\]
\[
+ \gamma\lambda\bigl[R_{t+2} + \gamma\hat{v}(S_{t+2},\mathbf{w}_{t+1}) - \hat{v}(S_{t+1},\mathbf{w}_t)\bigr]
\]
\[
+ (\gamma\lambda)^2\bigl[R_{t+3} + \gamma\hat{v}(S_{t+3},\mathbf{w}_{t+2}) - \hat{v}(S_{t+2},\mathbf{w}_{t+1})\bigr]
\]
\[
+ \cdots + (\gamma\lambda)^{k-1}\bigl[R_{t+k} + \gamma\hat{v}(S_{t+k},\mathbf{w}_{t+k-1}) - \hat{v}(S_{t+k-1},\mathbf{w}_{t+k-2})\bigr]
\]

\text{Notice the telescoping: the } $+\gamma\hat{v}(S_{i+1},\mathbf{w}_i)$ \text{ term from step } $i$ 
\text{ cancels with the } $-\hat{v}(S_{i+1}, \mathbf{w}_i)$ \text{ term from step } $i+1$ 
\text{ (after absorbing the } $\gamma\lambda$ \text{ factor). What remains is:}

\[
= -\hat{v}(S_t,\mathbf{w}_{t-1}) + \sum_{i=t}^{t+k-1}(\gamma\lambda)^{i-t}R_{i+1} 
+ (\gamma\lambda)^{k-1}\cdot\gamma\hat{v}(S_{t+k},\mathbf{w}_{t+k-1})
\]

\text{Rearranging:}
\[
\hat{v}(S_t,\mathbf{w}_{t-1}) + \sum_{i=t}^{t+k-1}(\gamma\lambda)^{i-t}\delta_i'
= \sum_{i=t}^{t+k-1}(\gamma\lambda)^{i-t}R_{i+1} + \gamma^k\lambda^{k-1}\hat{v}(S_{t+k},\mathbf{w}_{t+k-1})
\]

\text{Which matches the definition of } $G_{t:t+k}^\lambda$ \text{ — a } $k$\text{-step } 
$\lambda$\text{-return bootstrapping at } $S_{t+k}$\text{. } $\square$
\\\\
\textbf{Exercise 12.6} \quad Modify the pseudocode for Sarsa($\lambda$) to use dutch traces (12.11) without the other distinctive features of a true online algorithm. Assume linear function approximation and binary features.

\textbf{Solution:}
\begin{algorithm}
\caption{Sarsa($\lambda$) with Dutch traces and binary features}
\begin{algorithmic}[1]

\State \textbf{Input:} a function $\mathcal{F}(s, a)$ returning the set of (indices of) active features for $s, a$
\State \textbf{Input:} a policy $\pi$
\State \textbf{Algorithm parameters:} step size $\alpha > 0$, trace decay rate $\lambda \in [0,1]$, small $\varepsilon > 0$
\State \textbf{Initialize:} $\mathbf{w} = (w_1, \ldots, w_d)^\top \in \mathbb{R}^d$ (e.g., $\mathbf{w} = \mathbf{0}$), $\mathbf{z} = (z_1, \ldots, z_d)^\top \in \mathbb{R}^d$

\For{each episode}
    \State Initialize $S$
    \State Choose $A \sim \pi(\cdot|S)$ or $\varepsilon$-greedy according to $\hat{q}(S, \cdot, \mathbf{w})$
    \State $\mathbf{z} \leftarrow \mathbf{0}$
    \For{each step of episode}
        \State Take action $A$, observe $R, S'$
        \State $\delta \leftarrow R$
        \For{$i$ in $\mathcal{F}(S, A)$}
            \State $\delta \leftarrow \delta - w_i$
            \State $z_i \leftarrow (1 - \alpha z_i) z_i + 1$ \Comment{Dutch trace update}
        \EndFor
        \If{$S'$ is terminal}
            \State $\mathbf{w} \leftarrow \mathbf{w} + \alpha \delta \mathbf{z}$
            \State Go to next episode
        \EndIf
        \State Choose $A' \sim \pi(\cdot|S')$ or $\varepsilon$-greedy according to $\hat{q}(S', \cdot, \mathbf{w})$
        \For{$i$ in $\mathcal{F}(S', A')$}
            \State $\delta \leftarrow \delta + \gamma w_i$
        \EndFor
        \State $\mathbf{w} \leftarrow \mathbf{w} + \alpha \delta \mathbf{z}$
        \State $\mathbf{z} \leftarrow \gamma \lambda \mathbf{z}$
        \State $S \leftarrow S'$; $A \leftarrow A'$
    \EndFor
\EndFor

\end{algorithmic}
\end{algorithm}
\\\\
\textbf{Exercise 12.7:} Generalize the three recursive equations (12.18, 12.19, 12.20) 
to their truncated versions, defining $G_{t:h}^{\lambda s}$ \text{and} $G_{t:h}^{\lambda a}$.
\\
\\

\begin{align*}
G_{t:h} &=
\begin{cases}
R_{t+1} + \gamma G_{t+1:h}, & t < h \\
0, & t = h
\end{cases} \\[1em]
G_{t:h}^\lambda &=
\begin{cases}
R_{t+1} + \gamma \left[(1-\lambda)V(S_{t+1}) + \lambda G_{t+1:h}^\lambda\right], & t < h \\
V(S_h), & t = h
\end{cases} \\[1em]
v_{t:h}(S_t) &=
\mathbb{E}\left[
\begin{cases}
R_{t+1} + \gamma v_{t+1:h}(S_{t+1}), & t < h \\
0, & t = h
\end{cases}
\right]
\end{align*}
\\\\
\textbf{Exercise 12.8:} Prove that (12.24) becomes exact if the value function does not
change. To save writing, consider the case of $t = 0$, and use the notation
$V_k \doteq \hat{v}(S_k, \mathbf{w})$.

\bigskip

\textbf{Solution:} Begin by expanding $G_0^{\lambda s}$ using (12.22):

\begin{equation}
G_0^{\lambda s} = \rho_0\Big(R_1 + \gamma_1\big((1-\lambda_1)V_1
+ \lambda_1 G_1^{\lambda s}\big)\Big) + (1-\rho_0)V_0
\end{equation}

Recall the state-based TD error (12.23):

\begin{equation}
\delta_k^s \doteq R_{k+1} + \gamma_{k+1}V_{k+1} - V_k
\end{equation}

We can rearrange to get:

\begin{equation}
R_{k+1} + \gamma_{k+1}V_{k+1} = \delta_k^s + V_k
\end{equation}

Now substitute back into the expansion of $G_0^{\lambda s}$ and expand $G_1^{\lambda s}$
recursively. At each step $k$, the $(1 - \rho_k)V_k$ terms and the bootstrapped $V_k$
terms from $\delta_k^s$ telescope cleanly because the value function is fixed. After full
recursive expansion, all intermediate $V_k$ terms cancel except $V_0$, leaving:

\begin{equation}
G_0^{\lambda s} = V_0 + \rho_0 \sum_{k=0}^{\infty} \delta_k^s
\prod_{i=1}^{k} \gamma_i \lambda_i \rho_i
\end{equation}

which is exactly (12.24) for $t = 0$. The approximation in (12.24) arises only when the
value function changes between updates (i.e.\ $\mathbf{w}$ is updated), so that
$\hat{v}(S_k, \mathbf{w}_t) \neq \hat{v}(S_k, \mathbf{w}_k)$. When $\mathbf{w}$ is
fixed, the telescoping is exact. $\square$

\bigskip

\textbf{Exercise 12.9:} The truncated version of the general off-policy return is denoted
$G_{t:h}^{\lambda s}$. Guess the correct equation, based on (12.24).

\bigskip

\textbf{Solution:} By analogy with (12.24), we truncate the infinite sum at horizon $h$:

\begin{equation}
G_{t:h}^{\lambda s} \doteq \hat{v}(S_t, \mathbf{w}_t) + \rho_t \sum_{k=t}^{h-1}
\delta_k^s \prod_{i=t+1}^{k} \gamma_i \lambda_i \rho_i
\end{equation}

The structure is identical to (12.24) except the sum runs only to $h-1$ rather than
$\infty$. Bootstrapping beyond $h$ is handled implicitly through the $\delta_k^s$ terms,
which already contain $\hat{v}(S_{k+1}, \mathbf{w})$. As with (12.24), this becomes
exact when the value function does not change. $\square$
latex\documentclass{article}
\usepackage{amsmath}
\usepackage{amssymb}

\\\\
\section*{Exercise 12.10}

\subsection*{Question}

Prove that (12.27) becomes exact if the value function does not change. To save writing, 
consider the case of $t = 0$, and use the notation $Q_k = \hat{q}(S_k, A_k, \mathbf{w})$.

\subsection*{Answer}

We wish to show that the eligibility-trace-based update is exact when $\mathbf{w}$ is fixed. 
We expand $G_0^{\lambda a}$ term by term by unrolling the recursive definition.

The expected TD error at step $k$ is:

$$\delta_k^s = R_{k+1} + \gamma \bar{V}(S_{k+1}) - Q_k$$

where $\bar{V}(S_k) = \sum_a \pi(a|S_k)\hat{q}(S_k, a, \mathbf{w})$ is the expected value 
under the policy.

\textbf{At $k=0$:} The first term of the unrolled return is:

$$Q_0 + \delta_0^s = Q_0 + R_1 + \gamma \bar{V}(S_1) - Q_0 = R_1 + \gamma \bar{V}(S_1)$$

\textbf{At $k=1$:} Adding the next term $\gamma\lambda\delta_1^s$:

$$R_1 + \gamma \bar{V}(S_1) + \gamma\lambda\delta_1^s 
= R_1 + \gamma \bar{V}(S_1) + \gamma\lambda\bigl(R_2 + \gamma \bar{V}(S_2) - Q_1\bigr)$$

\textbf{At $k=2$:} Adding the next term $(\gamma\lambda)^2\delta_2^s$:

$$R_1 + \gamma \bar{V}(S_1) + \gamma\lambda\bigl(R_2 + \gamma \bar{V}(S_2) - Q_1\bigr) 
+ (\gamma\lambda)^2\bigl(R_3 + \gamma \bar{V}(S_3) - Q_2\bigr)$$

The general pattern is now clear. At each step $k \geq 1$, the term $-Q_k$ from 
$\delta_k^s$ and the term $\gamma \bar{V}(S_k)$ from $\delta_{k-1}^s$ appear together. 
When $\mathbf{w}$ is fixed, we have:

$$\bar{V}(S_k) = \sum_a \pi(a|S_k)\hat{q}(S_k, a, \mathbf{w})$$

and since the action $A_k$ is chosen according to $\pi$, the expected value satisfies 
$\gamma\bar{V}(S_k) = \gamma Q_k$, so these terms cancel exactly:

$$\gamma \bar{V}(S_k) - \gamma Q_k = 0 \quad \text{for all } k \geq 1$$

This telescoping cancellation propagates through all intermediate steps, leaving only 
the discounted reward stream and the terminal bootstrap term, which is precisely the 
definition of the $\lambda$-return $G_0^{\lambda a}$. Therefore (12.27) becomes exact 
when the value function does not change. \hfill$\square$
\\\\
\section*{Exercise 12.11}

\\section*{Exercise 12.11}

\subsection*{Question}

The truncated version of the general off-policy return is denoted $G_{t:h}^{\lambda a}$. 
Guess the correct equation for it, based on (12.27).

\subsection*{Answer}

By analogy with the state case, the truncated action-value $\lambda$-return is obtained 
by simply capping the summation at horizon $h$ instead of running to infinity. Recall 
the full off-policy return (12.27):

$$G_t^{\lambda a} \approx \hat{q}(S_t, A_t, \mathbf{w}_t) + \sum_{k=t}^{\infty} \delta_k^s 
\prod_{i=t+1}^{k} \gamma_i \lambda_i \rho_i$$

where the expected TD error is:

$$\delta_k^s = R_{k+1} + \gamma_{k+1}\bar{V}(S_{k+1}) - \hat{q}(S_k, A_k, \mathbf{w}_k)$$

and $\bar{V}(S_k) = \sum_a \pi(a|S_k)\hat{q}(S_k, a, \mathbf{w})$ is the expected value 
under the target policy, and $\rho_i = \frac{\pi(A_i|S_i)}{b(A_i|S_i)}$ are the 
importance sampling ratios for off-policy correction.

The truncated version is then:

$$G_{t:h}^{\lambda a} \doteq \hat{q}(S_t, A_t, \mathbf{w}_t) + \sum_{k=t}^{h-1} \delta_k^s 
\prod_{i=t+1}^{k} \gamma_i \lambda_i \rho_i$$

In the on-policy case $\rho_i = 1$ for all $i$, so the importance sampling ratios 
disappear and we recover the simpler on-policy form. \hfill$\square$
\\\\
\section*{Exercise 12.12}

\subsection*{Question}

Show in detail the steps outlined above for deriving (12.29) from (12.27). Start with 
the update (12.15), substitute $G_t^{\lambda a}$ from (12.26) for $G_t^{\lambda}$, then 
follow similar steps as led to (12.25).

\subsection*{Answer}

We begin with the semi-gradient update (12.15) for action values:

$$\mathbf{w}_{t+1} \doteq \mathbf{w}_t + \alpha\left[G_t^{\lambda a} - 
\hat{q}(S_t, A_t, \mathbf{w}_t)\right]\nabla\hat{q}(S_t, A_t, \mathbf{w}_t)$$

From (12.27), we have:

$$G_t^{\lambda a} - \hat{q}(S_t, A_t, \mathbf{w}_t) = \sum_{k=t}^{\infty} \delta_k^s 
\prod_{i=t+1}^{k} \gamma_i \lambda_i \rho_i$$

where the expected action-based TD error is:

$$\delta_k^s = R_{k+1} + \gamma_{k+1}\bar{V}(S_{k+1}) - \hat{q}(S_k, A_k, \mathbf{w}_k)$$

Substituting into the update and summing over all timesteps:

$$\sum_{t=0}^{\infty}(\mathbf{w}_{t+1} - \mathbf{w}_t) \approx \sum_{t=0}^{\infty} 
\sum_{k=t}^{\infty} \alpha \nabla\hat{q}(S_t, A_t, \mathbf{w}_t) \delta_k^s 
\prod_{i=t+1}^{k} \gamma_i \lambda_i \rho_i$$

Swapping the order of summation using the rule 
$\sum_{t=x}^{\infty}\sum_{k=t}^{\infty} = \sum_{k=x}^{\infty}\sum_{t=x}^{k}$:

$$= \sum_{k=0}^{\infty} \alpha \delta_k^s \sum_{t=0}^{k} \nabla\hat{q}(S_t, A_t, 
\mathbf{w}_t) \prod_{i=t+1}^{k} \gamma_i \lambda_i \rho_i$$

This is in the form of a backward-view TD update if the inner sum can be written and 
updated incrementally as an eligibility trace. We now show this is possible. Define:

$$\mathbf{z}_k \doteq \sum_{t=0}^{k} \nabla\hat{q}(S_t, A_t, \mathbf{w}_t) 
\prod_{i=t+1}^{k} \gamma_i \lambda_i \rho_i$$

We can verify this can be updated incrementally from $\mathbf{z}_{k-1}$:

\begin{align*}
\mathbf{z}_k &= \sum_{t=0}^{k} \nabla\hat{q}(S_t, A_t, \mathbf{w}_t) 
\prod_{i=t+1}^{k} \gamma_i \lambda_i \rho_i \\
&= \sum_{t=0}^{k-1} \nabla\hat{q}(S_t, A_t, \mathbf{w}_t) 
\prod_{i=t+1}^{k} \gamma_i \lambda_i \rho_i + \nabla\hat{q}(S_k, A_k, \mathbf{w}_k) \\
&= \gamma_k \lambda_k \rho_k \sum_{t=0}^{k-1} \nabla\hat{q}(S_t, A_t, \mathbf{w}_t) 
\prod_{i=t+1}^{k-1} \gamma_i \lambda_i \rho_i + \nabla\hat{q}(S_k, A_k, \mathbf{w}_k) \\
&= \gamma_k \lambda_k \rho_k \mathbf{z}_{k-1} + \nabla\hat{q}(S_k, A_k, \mathbf{w}_k)
\end{align*}

Changing the index from $k$ to $t$, this is exactly the eligibility trace update (12.29):

$$\boxed{\mathbf{z}_t \doteq \gamma_t \lambda_t \rho_t \mathbf{z}_{t-1} + 
\nabla\hat{q}(S_t, A_t, \mathbf{w}_t)}$$

which is the general accumulating trace update for action values, directly analogous 
to the state-value trace (12.25), with $\nabla\hat{v}(S_t, \mathbf{w}_t)$ replaced by 
$\nabla\hat{q}(S_t, A_t, \mathbf{w}_t)$ throughout. \hfill$\square$
\\\\
\section*{Exercise 12.13}

\subsection*{Question}

What are the dutch-trace and replacing-trace versions of off-policy eligibility traces 
for state-value and action-value methods?

\subsection*{Answer}

We take the general accumulating trace updates (12.25) and (12.29) and substitute in 
the Dutch trace and replacing trace definitions.

\subsubsection*{State-Value Methods}

The general accumulating trace update (12.25) is:

$$\mathbf{z}_t \doteq \rho_t(\gamma_t\lambda_t\mathbf{z}_{t-1} + \nabla\hat{v}(S_t, \mathbf{w}_t))$$

\textbf{Replacing trace} (state-value):

$$z_{t,i} \doteq \rho_t\left(\gamma_t\lambda_t z_{t-1,i} + \begin{cases} 1 & \text{if } x_{t,i} = 1 \\ 0 & \text{otherwise}\end{cases}\right)$$

where $x_{t,i}$ is the $i$-th component of the feature vector at time $t$.

\textbf{Dutch trace} (state-value):

$$\mathbf{z}_t \doteq \rho_t\left(\gamma_t\lambda_t\mathbf{z}_{t-1} + 
(1 - \alpha\gamma_t\lambda_t\rho_t\mathbf{z}_{t-1}^\top\nabla\hat{v}(S_t, \mathbf{w}_t))
\nabla\hat{v}(S_t, \mathbf{w}_t)\right)$$

\subsubsection*{Action-Value Methods}

The general accumulating trace update (12.29) is:

$$\mathbf{z}_t \doteq \gamma_t\lambda_t\rho_t\mathbf{z}_{t-1} + \nabla\hat{q}(S_t, A_t, \mathbf{w}_t)$$

\textbf{Replacing trace} (action-value):

$$z_{t,i} \doteq \gamma_t\lambda_t\rho_t z_{t-1,i} + \begin{cases} 1 & \text{if } x_{t,i} = 1 \\ 0 & \text{otherwise}\end{cases}$$

\textbf{Dutch trace} (action-value):

$$\mathbf{z}_t \doteq \gamma_t\lambda_t\rho_t\mathbf{z}_{t-1} + 
(1 - \alpha\gamma_t\lambda_t\rho_t\mathbf{z}_{t-1}^\top\nabla\hat{q}(S_t, A_t, \mathbf{w}_t))
\nabla\hat{q}(S_t, A_t, \mathbf{w}_t)$$

In all cases the substitution is straightforward: the Dutch and replacing trace variants 
modify how the current gradient term is added to the decayed previous trace, while the 
off-policy correction via $\rho_t$ carries through identically from the accumulating 
case. For action-value methods, $\nabla\hat{v}(S_t, \mathbf{w}_t)$ is replaced 
throughout by $\nabla\hat{q}(S_t, A_t, \mathbf{w}_t)$. \hfill$\square$
\\\\


\end{document}










